%% ma: L10-B15-0039
%% lop: 10
%% chuong: 6
%% bai: 15
%% dang: 10.15.2
%% loai: TL
%% muc: VD
%% dap_an: "-"
%% nguon: "Ảnh giáo viên gửi 21/11/2026: Chuyên đề VI – Toán 10 – Chương VI (Hàm số – đồ thị và ứng dụng), tr. 199–203, Dạng 2 Câu 8a"
%% kien_thuc: "Lớp 10 Bài 15 (tập xác định); hằng đẳng thức (lớp 8)"
%% trang_thai: chinh-thuc
%% ghi_chu: "Dạng do giáo viên cung cấp (21/11/2026). Tách từ Câu 8 (ý a)."
\begin{ex}
Tìm $m$ để hàm số $y=\dfrac{2x+1}{\sqrt{x^2-6x+m-2}}$ xác định trên $\mathbb R$.
\loigiai{Cần $x^2-6x+m-2>0$ với mọi $x$. Mà $x^2-6x+m-2=(x-3)^2+m-11$ có giá trị nhỏ nhất $m-11$ (tại $x=3$). Vậy $m-11>0\Leftrightarrow m>11$.}
\end{ex}
